\documentclass[
  aps,
  prb,
  reprint,
  amsmath,
  amssymb,
  superscriptaddress,
  longbibliography,
  nofootinbib,
  floatfix
]{revtex4-2}

\usepackage[T1]{fontenc}
\usepackage{lmodern}
\usepackage{graphicx}
\usepackage{bm}
\usepackage{microtype}
\usepackage[hidelinks]{hyperref}

\graphicspath{{figures/}}

% 放宽跨栏浮动图的排版限制。
\setcounter{topnumber}{4}
\setcounter{dbltopnumber}{4}
\setcounter{totalnumber}{6}
\renewcommand{\topfraction}{0.9}
\renewcommand{\dbltopfraction}{0.9}
\renewcommand{\bottomfraction}{0.5}
\renewcommand{\textfraction}{0.07}
\renewcommand{\floatpagefraction}{0.7}
\renewcommand{\dblfloatpagefraction}{0.7}

\newcommand{\Heff}{H_{\mathrm{eff}}}
\newcommand{\Jeff}{J_{\mathrm{eff}}}
\newcommand{\ceff}{c_{\mathrm{eff}}}
\newcommand{\qpeak}{q_{\mathrm{peak}}}
\newcommand{\Szz}{S^{zz}(q,\omega)}
\newcommand{\avg}[1]{\left\langle #1\right\rangle}
\newcommand{\pentamer}{pentamer}

\begin{document}

\title{Local multiplets, effective-chain dynamics, and finite-size momentum
reconstruction in a spin-\texorpdfstring{$1/2$}{1/2} Heisenberg pentamer chain}

\author{Author Name}
\email{corresponding.author@example.edu}
\affiliation{Department and Institution, City, Postal Code, Country}

\date{\today}

\begin{abstract}
In chains of spin clusters, the symmetry of a single cluster can decide which
excitations a scattering experiment sees.
We establish this for the spin-$1/2$ Heisenberg chain with a five-site unit cell.
Because the cell holds an odd number of spins, its ground state for positive
couplings is not a singlet but a spin-$1/2$ doublet that is even under reflection.
The inter-cell bond couples these doublets into an antiferromagnetic effective
Heisenberg chain, with corrections quadratic in the inter-cell coupling, and
entanglement, finite-size gaps, correlations, and a dominant $q=\pi$ response
identify a $c=1$ critical chain at the couplings studied.
Reflection parity makes the lowest excited doublet invisible to a $q=\pi$ spin
probe, so spin-$3/2$ multiplets carry $84\%$--$98\%$ of the local inelastic
weight.
Complete connected real-time spectra at $J_2/J_1=\pm0.5$, $J_3/J_1=0.5$, and
$L=40$ compare the low- and high-energy responses at a common finite-time
resolution, without identifying individual bands with local multiplets.
For negative intra-cell exchange the cell ground state becomes a spin-$3/2$
quartet whose shape factor favors the $0.60\pi$ replica among the folded
antiferromagnetic momenta, and finite chains move the dominant structure-factor
momentum from $0.60\pi$ toward $\pi$.
Because the momentum and finite-size-flow diagnostics of this reconstruction
select different coupling intervals, the present data establish a finite-size
reconstruction without resolving a separate thermodynamic phase.
\end{abstract}

\maketitle

\section{Introduction}
\label{sec:introduction}

Clusterized quantum magnets provide a direct route from the spectrum of a
few-spin unit to the collective behavior of a quantum many-body system.
Their local multiplets can survive as low-energy degrees of freedom, become
dispersive, or be removed from an experimental response by symmetry.
Clusterized chains therefore separate three questions that are often mixed
together in a conventional phase diagram: which local state remains active, how
neighboring units couple, and which resulting excitations are visible.

The spin-$1/2$ tetramer chain is a well-studied reference of this type, because
spin order, entanglement structure, and excitation spectra have been organized
within a common coupling diagram~\cite{li2026tetramer,gong2008tetramer}.
Experimentally, the tetramer compound CuInVO$_5$ has been characterized as a
multimerized magnet with a singlet unit-cell ground state, and its magnetic
spectrum has been discussed in terms of both local multiplet and spinon
excitations~\cite{hase2016cuinvo5,reja2019cuinvo5,singhania2018}.
Trimer chains, by contrast, host a doublet unit cell and support fractionalized
and composite excitations that have been resolved in
RIXS~\cite{cheng2022trimer,cheng2024trimerfield,li2025rixstrimer,prabhakar2025trimer,bera2022trimer}.
Beyond the spectrum, the dimerized and tetramerized limits are classified by
nonlocal string order and symmetry-protected
topology~\cite{guwen2009,pollmann2010,pollmann2012,han2017bcfg,ma2023aspt},
diagnosed by order parameters and entanglement measures that do not rely on a
local multiplet picture~\cite{hung2015,lu2020,scheie2021}.
Related one-dimensional models further show that an intermediate deconfined
critical state can separate two gapped cluster
phases~\cite{jiang2019dqcp,roberts2019dqcp,huang2019dqcp,yang2021dqcp}, a
possibility that has to be tested explicitly for an odd unit cell.
Even in the molecular limit, the collective spin dynamics of a closed cluster is
not simply a sum of local transitions~\cite{ummethum2012}.

Replacing four spins by five is not a minor geometric modification.
A unit cell with an odd number of half-integer spins retains a half-integer
degree of freedom unless symmetry breaking or ground-state degeneracy
intervenes.
Lieb--Schultz--Mattis-type results therefore rule out a unique
symmetry-preserving gapped ground state under the relevant translation and
spin-rotation assumptions~\cite{lieb1961,ogata2019,affleck1987}, in close analogy
with the Haldane conjecture for integer-spin chains~\cite{haldane1983} and the
dimerized-chain crossovers studied earlier~\cite{hida1992,chitra1995}.
The odd cell thus changes the question.
It is no longer which gapped cluster phase forms, but which cell state carries the
low-energy physics, whether the resulting chain is the natural $c=1$ Luttinger
liquid~\cite{giamarchi2003,calabrese2004}, and which of its excitations a spin
probe can reach.

Here we study a linear Heisenberg pentamer chain with two intra-cell exchanges,
$J_1$ and $J_2$, and an inter-cell exchange $J_3$.
We set $J_1=1$ and define $\alpha=J_2/J_1$ and $\beta=J_3/J_1$.
The geometry differs from previously studied mixed-spin or differently connected
pentamers~\cite{haraldsen2011pentamer}.
The isolated cell is the natural starting point: its ground multiplet fixes the
effective degree of freedom of the chain, and its reflection parity fixes which
local excitations a spin probe can excite.
We therefore determine the complete isolated-cell spectrum, construct the
projected low-energy chain, and test this picture against static and dynamical
many-body calculations.

Three results organize this work.
First, throughout the sampled positive-$\alpha$ region, the pentamer ground
state is an even-parity spin-$1/2$ doublet and positive $\beta$ generates an
antiferromagnetic effective chain.
Second, reflection parity makes the lowest excited doublet essentially dark at
$q=\pi$ and concentrates the local inelastic weight in spin-$3/2$ multiplets, in
parallel with the momentum-selective response discussed for related clusterized
chains~\cite{gao2024srv3o9,schlappa2018,ishii2025kedge}.
Third, negative $\alpha$ switches the isolated cell to a spin-$3/2$ quartet whose
shape factor favors $0.60\pi$ among the three inequivalent physical momenta
folded from effective-chain $Q=\pi$.
Finite chains develop a momentum reconstruction from that value toward $\pi$.
Two independent diagnostics of this reconstruction select different coupling
intervals, leaving a finite-size crossover interpretation compatible with the
data but no resolved thermodynamic phase boundary.
Complete connected spectra at $(\alpha,\beta)=(\pm0.5,0.5)$ supplement this
static comparison at fixed $L=40$ and a common finite-time resolution.
These results rest on the standard methodology of one-dimensional magnetism: matrix-product-state
DMRG~\cite{white1992,schollwoeck2005,schollwoeck2011} combined with
correction-vector and real-time spectral
techniques~\cite{kuehner1999,jeckelmann2002,vidal2004,paeckel2019}.
Extensions of these methods to finite temperature and entanglement-based
diagnostics provide related approaches~\cite{barthel2009,fishman2022,keselman2020}.

Section~\ref{sec:methods} defines the model, the projection, and the numerical
criteria.
Section~\ref{sec:results} follows the positive-coupling chain from the isolated
cell to the collective spectrum and then turns to negative $\alpha$.
Section~\ref{sec:discussion} compares the results with the tetramer chain and
states their consequences for scattering experiments.

\section{Model and numerical methods}
\label{sec:methods}

\subsection{Hamiltonian and isolated-cell projection}

The Hamiltonian of the pentamer chain is
\begin{align}
H={}&J_1\sum_n\left(\bm S_{n,2}\!\cdot\!\bm S_{n,3}
 +\bm S_{n,3}\!\cdot\!\bm S_{n,4}\right) \notag\\
&+J_2\sum_n\left(\bm S_{n,1}\!\cdot\!\bm S_{n,2}
 +\bm S_{n,4}\!\cdot\!\bm S_{n,5}\right) \notag\\
&+J_3\sum_n\bm S_{n,5}\!\cdot\!\bm S_{n+1,1},
\label{eq:hamiltonian}
\end{align}
where $n$ labels the five-site unit cells and $\bm S_{n,a}$ is the spin-$1/2$
operator at site $a$ of cell $n$.
The intra-cell bonds are the two outermost bonds at $J_2$ and the two central
bonds at $J_1$; a single bond $J_3$ connects neighboring cells.
We work in units of $J_1$ and set $\alpha=J_2/J_1$ and $\beta=J_3/J_1$ throughout.

The isolated pentamer has a 32-dimensional Hilbert space.
We diagonalize it completely and resolve each degenerate energy subspace by total
spin $S$ and by reflection parity about the cell center.
Spectral weights are evaluated with projectors onto complete multiplets rather
than with individual eigenvectors, whose orientation inside a degenerate subspace
is arbitrary.
This matters here because the multiplet structure, not merely the energy
ordering, controls the dynamical selection rules derived below.

If $P$ projects onto the ground-state doublet, the local spin operators act
within that subspace as
\begin{equation}
P\bm S_{n,a}P=c_a\bm T_n,
\label{eq:projection}
\end{equation}
where $\bm T_n$ is an effective spin-$1/2$ operator and the $c_a$ are
cell-dependent projection coefficients.
To leading order in the inter-cell coupling the chain is then described by
\begin{equation}
\Heff=NE_5+\Jeff\sum_n\bm T_n\!\cdot\!\bm T_{n+1}+O(\beta^2),
\qquad \Jeff=\beta c_1c_5,
\label{eq:effective}
\end{equation}
with $N$ the number of cells and $E_5$ the isolated-cell ground-state energy.
Only the end sites $a=1$ and $5$ enter $\Jeff$, because they are the sites touched
by the inter-cell bond.
Reflection about the cell center maps site 1 onto site 5.
Within an isolated ground multiplet of definite reflection parity, the parity
phases of bra and ket cancel, giving $c_5=c_1$ for either parity.
Thus $\Jeff=\beta c_1^2$ is antiferromagnetic for $\beta>0$ and nonzero endpoint
coefficient; its magnitude measures how much of the cell spin resides on the ends.
Even parity is needed for the selection rule below, not for this sign argument.
The same coefficients determine what a scattering probe sees.
Because each effective spin is spread over the cell with amplitudes $c_a$, a
probe at physical momentum $q$ couples to it with the cell shape factor
\begin{equation}
F(q)=\Bigl|\sum_{a=1}^{5}c_a e^{iqa}\Bigr|^2,
\label{eq:shape}
\end{equation}
which weights the folded effective-chain response in the physical Brillouin zone.
With physical site spacing set to one, the cell momentum is
$Q=5q\pmod{2\pi}$.
For spin operators normalized by $1/\sqrt{L}$ in the physical chain and
$1/\sqrt{N}$ in the effective chain, the leading projected response is
\begin{equation}
S^{zz}_{\mathrm{phys}}(q,\omega)\simeq
\frac{F(q)}{5}S^{zz}_{\mathrm{eff}}(Q,\omega).
\label{eq:folded-response}
\end{equation}
For a quartet ground multiplet the same relation uses effective spin-$3/2$
operators in place of $\bm T_n$.
The shape factor therefore ranks replicas at a fixed $Q$; its maximum alone
does not determine the maximum of the full chain response.
We test the projection in two ways.
Exact diagonalization of two coupled pentamers measures the departure of the
exact ground-state energy from the projected form, and finite-chain calculations
with $N=8,12,16$ cells test whether the same leading power persists in the
extended system.
Figure~\ref{fig:model} summarizes the geometry, the isolated-cell multiplet
structure, and the projection that defines the effective chain.

\begin{figure*}[t]
  \includegraphics[width=\textwidth]{fig1_model_projection.png}
  \caption{Model and low-energy projection.
  The five-site unit cell defines the couplings $J_1$, $J_2$, and $J_3$.
  Isolated-pentamer multiplets are resolved by total spin and reflection parity,
  and the even-parity ground-state doublet is projected onto an effective
  spin-$1/2$ degree of freedom that mediates the inter-cell exchange.}
  \label{fig:model}
\end{figure*}

\subsection{Finite-chain calculations and convergence criteria}
\label{sec:convergence}

Finite chains use open boundaries and $L=5N$ physical spins.
Ground states are obtained in fixed total-$S^z$ sectors with density-matrix
renormalization group calculations in a matrix-product-state
representation~\cite{white1992,schollwoeck2005,schollwoeck2011}.
Sector-specific initial states are used rather than assuming that a single
initialization samples every low-energy solution.

The positive-region data set contains 105 converged calculations at
$L=40,60,80$.
For negative $\alpha$, we initialize the $S^z=0$ sector from one of two distinct
N\'eel-type product states, and the $S^z=1$ sector from a single fixed product
state.
A calculation is retained only if the variance per site satisfies
$\mathrm{Var}(H)/L\le5\times10^{-8}$, the final truncation error satisfies
$\epsilon_{\mathrm{trunc}}\le10^{-9}$, the measured magnetization satisfies
$|S^z-S^z_{\mathrm{target}}|\le10^{-8}$, and the sector gap
$E_0(S^z=1)-E_0(S^z=0)$ is no smaller than $-10^{-8}$, so that the $S^z=0$ sector
holds the ground state; each calculation uses at most 64 sweeps.
The two $S^z=0$ initializations probe sensitivity to the starting state.

Every calculation was re-evaluated before use, and only calculations that pass
these criteria enter the analysis below.
These criteria test the variational state at a given solver setting; the
sensitivity to the eigensolver itself is examined separately in
Sec.~\ref{sec:krylov}.

\subsection{Entanglement, correlation, and momentum diagnostics}
\label{sec:diagnostics}

For an open critical chain the bond entropy is compared with the conformal form
\begin{equation}
S(b)=a_\nu+\frac{\ceff}{6}
\ln\left[\frac{2L}{\pi}\sin\left(\frac{\pi b}{L}\right)\right],
\label{eq:entropy}
\end{equation}
where the intercept $a_\nu$ depends on the bond class $\nu=b\bmod 5$ or $b\bmod 10$
and a common slope defines the descriptive coefficient $\ceff$.
Fits use central windows $fL\le b\le(1-f)L$ with $f=0.10$, 0.15, 0.20, and 0.25.
Equation~\eqref{eq:entropy} is the standard conformal result for the entanglement
entropy of an open critical
chain~\cite{calabrese2004,giamarchi2003}; because the five-site unit cell makes
bond positions inequivalent, we treat the intercept as bond-class dependent and
refer to the common slope as the effective entropy coefficient.

Correlation envelopes, finite-size sector gaps, the static structure factor, and
connected bond--bond correlations provide complementary diagnostics.
Momentum-resolved information is obtained from the static structure factor
$S(q)=\sum_{r}e^{iqr}\avg{\bm S_0\!\cdot\!\bm S_r}$ and from the dynamical
structure factor defined in Sec.~\ref{sec:dsf}.

The negative-$\alpha$ scan tracks the momentum $\qpeak$ at which $S(q)$ is
largest and the sharpness of that peak,
\begin{equation}
R_L=1-\frac{S(\qpeak-\delta)+S(\qpeak+\delta)}{2S(\qpeak)},
\label{eq:ratio}
\end{equation}
where $\delta$ is the momentum resolution $2\pi$ divided by the number of sites
in the central window.
A peak that sharpens with length, $\Delta R=R_{80}-R_{60}>0$, signals a developing
singularity, whereas a peak that broadens points away from one.
A candidate boundary is retained only when an adjacent coupling interval shows
both a momentum reconstruction and a sign change of $\Delta R$.
A momentum change without the corresponding $\Delta R$ condition is reported as
a finite-size reconstruction, not as a phase boundary.

\subsection{Dynamical structure factor}
\label{sec:dsf}

For the full-momentum spectra we use the connected operator
\begin{equation}
O_q=\frac{1}{\sqrt{L}}\sum_{j=1}^{L}e^{iqj}
\left(S_j^z-\avg{0|S_j^z|0}\right)
\label{eq:momentum-source}
\end{equation}
and the zero-temperature spectral representation
\begin{equation}
S^{zz}(q,\omega)=\sum_{m}\left|\avg{m|O_q|0}\right|^2
\delta\!\left(\omega-E_m+E_0\right)
\label{eq:dsf}
\end{equation}
with $\hbar=1$.
Real-time matrix-product-state evolution~\cite{vidal2004,paeckel2019} of
$|b_q\rangle=O_q|0\rangle$ gives
$C_q(t)=\langle b_q|e^{-i(H-E_0)t}|b_q\rangle$, including $t=0$.
The plotted finite-time estimator is
\begin{equation}
S_T^{zz}(q,\omega)=\frac{1}{\pi}\operatorname{Re}
\int_0^T dt\,e^{i\omega t}\left(1-\frac{t}{T}\right)C_q(t).
\label{eq:finite-time-dsf}
\end{equation}
The triangular time window sets the energy resolution and is part of the
spectral estimator.
No Gaussian smoothing, spatial window, interpolation, clipping, or momentum
reflection averaging is applied to these spectra.

Figure~\ref{fig:fullq} uses open chains with $L=40$ (eight complete cells),
$(\alpha,\beta)=(\pm0.5,0.5)$, and 41 independently evolved sources at
$q_m=2\pi m/40$, $m=0,\ldots,40$.
Second-order odd--even time evolution uses $TJ_1=60$, $\Delta tJ_1=0.05$,
maximum bond dimension $\chi=192$, and cutoff $10^{-10}$.
Five momenta, $q/\pi=0.4,0.7,1.0,1.3,1.6$, are recomputed with the same frozen
ground state at $\Delta tJ_1=0.025$ and $\chi=256$.
This joint time-step and bond-dimension check tests those five sources, rather
than establishing convergence at every momentum or separating the two errors.
The triangular-window line kernel has full width at half maximum
$5.56623/T=0.09277J_1$; a finer frequency display grid does not reduce this width.
For comparison with this scale, separate sector DMRG calculations measure
$E_0(S^z=1)-E_0(S^z=0)$ at the two plotted points, with joint budgets
$(\chi,\epsilon_{\mathrm{cut}})=(512,10^{-12})$ and $(768,10^{-13})$.
These are open-chain level differences, not thermodynamic bulk gaps; agreement
between the two budgets is not a rigorous error bound.
Both panels use one global color normalization,
$u(s)=\operatorname{asinh}(s/0.02)/\operatorname{asinh}(6/0.02)$ with
$s=J_1S_T^{zz}$, applied uniformly in $q$ and $\omega$.
This display mapping preserves the spectral data but enhances weak weight,
including any window sidelobes.

An earlier, separate $q=\pi$ comparison uses correction-vector
DMRG~\cite{kuehner1999,jeckelmann2002,nocera2016,nocera2018rixs}, which solves
the resolvent at $\omega+i\eta$, and a center-source, spatial-window
approximation to the time-dependent response.
The correction-vector implementation reproduces the complete $L=10$ exact
spectrum with relative $L_2$ error $9\times10^{-12}$ under its benchmark
normalization and broadening.
The finite-chain comparison in Fig.~\ref{fig:cv-tebd} is restricted to
$(\alpha,\beta)=(0.23,0.5)$ and is not an independent validation of the
complete-momentum spectra at $\alpha=\pm0.5$.

\subsection{Krylov-subspace stability control}
\label{sec:krylov}

At $(\alpha,\beta)=(-0.10,0.25)$ an earlier $L=80$ analysis gave an entropy
coefficient well below the values near unity found elsewhere
(Sec.~\ref{sec:saved-state}).
Because this point was also the most sensitive to the eigensolver configuration,
we performed a final $L=80$ control in which all physical and DMRG settings were held fixed while the
Krylov dimension was increased from $k=8$ to $k=12$.
Both initializations used 64 sweeps, a configured maximum bond dimension of 1200,
a cutoff of $10^{-11}$, and the same 401-point momentum grid.
Here $k$ is the Krylov-subspace dimension, distinct from the MPS bond dimension
$\chi$.

The archived closeout summary reports that the $k=12$ calculation satisfied
the same convergence criteria listed in Sec.~\ref{sec:convergence}.
Its variance per site was $1.9\times10^{-11}$, and the two sector gaps were
$7.2772\times10^{-5}$, agreeing to five significant figures.
Across 172 matched observables, the $k=8\rightarrow12$ differences were bounded by
$7.5\times10^{-12}$ for the spin gap, $8.4\times10^{-8}$ for the dominant
ratio, and $5.7\times10^{-8}$ for the entropy, while $\qpeak/\pi=1$ was
unchanged.
The largest retained link dimension was $264$--$265$ in the $S^z=0$ sector and
$186$--$190$ in $S^z=1$, nearly insensitive to the solver change and well below
the configured maximum of 1200, so truncation at this point is set by the
$10^{-11}$ cutoff rather than by the bond-dimension cap.
The terminal comparison files referenced by that summary are missing from the
current local source archive and must be restored for a source-level recheck
before release; the quoted $k=12$ values are historical summary records.

The $L=40$ and $L=60$ members of the negative-$\alpha$ size sequence in
Sec.~\ref{sec:saved-state} were obtained with the earlier setting $k=3$; the
quantitative statements there rest on the controlled $L=80$ endpoint.

\section{Results}
\label{sec:results}

\subsection{Low-energy spin doublet of the odd pentamer}
\label{sec:doublet}

At all 100 sampled points with $0<\alpha\le1$, the isolated-pentamer ground state
is an even-reflection-parity spin-$1/2$ doublet.
The lowest excited doublet is odd under reflection, whereas the lowest
spin-$3/2$ multiplet is even.
The five-spin unit therefore retains an active half-integer degree of freedom
instead of forming a fully quenched singlet cluster, in contrast to the
integer-spin and singlet cell sectors that organize the tetramer
problem~\cite{li2026tetramer,gong2008tetramer}.
How this cell-level doublet shapes the extended chain is the subject of the
following subsections.

At $\alpha=0.14$, 0.23, and 0.33, the projected coupling of
Eq.~\eqref{eq:effective} satisfies $\Jeff/\beta=c_1^2=0.440$, 0.432, and 0.419,
respectively.
The equal endpoint projection coefficients and positive $\beta$ generate an
antiferromagnetic effective spin-$1/2$ chain.
The residual between the projected and exact ground-state energies scales with
powers $1.97$--$1.99$ for two coupled pentamers and $1.948$--$1.950$ for
complete chains with $N=8,12,16$ cells.
The leading correction is thus the $O(\beta^2)$ term of
Eq.~\eqref{eq:effective}, and the agreement between the two-pentamer and
finite-chain exponents shows that it persists beyond two coupled cells.
The level flow and local spectral weights across the positive
$\alpha$ range are shown in Fig.~\ref{fig:multiplets}; the parity assignments
are described above.

\begin{figure*}[t]
  \includegraphics[width=\textwidth]{fig2_multiplet_parity.png}
  \caption{Multiplet flow and the $q=\pi$ parity filter for positive $\alpha$.
  Panel~rows show the isolated-pentamer levels resolved by total spin
  together with their projector-resolved $q=\pi$ spectral weights.
  Reflection-parity assignments are given in the text.
  The lowest excited doublet is dark, whereas even-parity spin-$3/2$ multiplets
  dominate the local inelastic response.}
  \label{fig:multiplets}
\end{figure*}

\subsection{A single resolved effective-chain regime at positive couplings}
\label{sec:positive-regime}

The strict positive-region mesh contains 105 converged calculations.
We test for a boundary with the sharpness of the $q=\pi$ structure-factor peak,
$R=1-S(\pi-\delta)/S(\pi)$, the one-sided form of Eq.~\eqref{eq:ratio} at
$\qpeak=\pi$.
A boundary along a slice of fixed $\beta$ would appear as a crossing of the
$L=40$, 60, and 80 curves of $R(\alpha)$, that is, as a sign change of their size
differences.
No such crossing occurs on any of the seven $\beta$ slices: the interior size
differences keep one sign at the $2\sigma$ level, also when any single
calculation is left out and when the least-converged points are examined
separately.
The data thus resolve no internal boundary at the sampled points in
$0<\alpha,\beta\le1$ at the mesh resolution.
This statement excludes the decoupled axes $\alpha=0$ and $\beta=0$.

The absence of a boundary leaves open what the single regime is; the low-energy
observables answer this.
The entanglement entropy measures the central charge, but in a five-site cell its
absolute $\ln L$ dependence carries boundary and unit-cell corrections.
We therefore fit the pentamer chain and a uniform spin-$1/2$ Heisenberg chain,
which has $c=1$, on the same open lengths and with the same window, and quote the
ratio of the two fitted slopes, in which these corrections largely cancel.
At $\alpha=0.23$ and weak inter-cell coupling, this matched-size ratio gives
$\ceff=1.02$--$1.03$ for open chains of $N=8$--$24$ pentamers ($L=40$--$120$), and
the bond entropy is stable against a doubling of the retained bond dimension (largest
relative change $2\times10^{-4}$).
Finite-size spin gaps extrapolate close to zero, the correlation envelope has
exponents $0.95$--$0.98$, and the static structure factor is dominated by
$q=\pi$.
Together these observables identify a gapless antiferromagnetic effective chain
at the tested points~\cite{descloizeaux1962,calabrese2004}.
The exploratory phase map and the converged finite-size mesh are contrasted in
Fig.~\ref{fig:positive-static}, and the entanglement calibration is detailed in
Fig.~\ref{fig:entanglement}.

\begin{figure*}[t]
  \includegraphics[width=0.98\textwidth]{fig3_positive_static.png}
  \caption{Positive-region ground-state diagnostics.
  The dense exploratory map is shown separately from the strict 105-calculation
  $L=40,60,80$ mesh that supplies the crossing test.
  No stable internal bracket is resolved by these data.}
  \label{fig:positive-static}
\end{figure*}

\begin{figure*}[t]
  \includegraphics[width=0.94\textwidth]{fig4_entanglement.png}
  \caption{Entanglement and finite-size low-energy diagnostics.
  Five inequivalent bond classes are calibrated against a uniform $c=1$
  Heisenberg chain of the same length.}
  \label{fig:entanglement}
\end{figure*}

\subsection{Reflection parity and the visible local excitation}
\label{sec:parity}

The isolated cell already determines which excitations a $q=\pi$ spin probe can
reach.
Within one cell, $S^z_{q=\pi}$ weights the five sites with alternating signs, a
pattern symmetric under the reflection $a\to6-a$; the operator is therefore even
and cannot connect the even ground doublet to an odd state.
The lowest excited doublet is odd, so it is not the lowest visible excitation:
its projector-resolved weight is below $10^{-25}$ at the three representative
positive-$\alpha$ points, a symmetry selection rule rather than a small numerical
residue.
The lowest even-parity spin-$3/2$ multiplet carries weight $0.33$, and
all spin-$3/2$ multiplets together account for $84\%$--$98\%$ of the
isolated-pentamer inelastic response.
The result identifies which local sector can contribute when collective states
form, and it provides a concrete spectral analogue of the
momentum-selective weight reported for related cluster
magnets~\cite{gao2024srv3o9,schlappa2018}.

Once cells couple, the question is whether these bright local states fall inside
the low-energy continuum of the effective chain, where they would mix with it, or
remain above it.
In the weak-coupling regime the effective-chain two-spinon continuum is bounded by the des Cloizeaux--Pearson lower edge
$\omega_{\mathrm L}(Q)=\pi\Jeff|\sin Q|/2$~\cite{descloizeaux1962} and the
standard upper edge $\omega_{\mathrm U}(Q)=\pi\Jeff|\sin(Q/2)|$,
where $Q=5q\pmod{2\pi}$ is the cell momentum.
At physical $q=\pi$ one also has $Q=\pi$.
Setting the upper edge there equal to the lowest local quartet gap defines a
formal embedding threshold at $\beta=0.104$, $0.176$, and $0.265$ for
$\alpha=0.14$, 0.23, and 0.33.
Measured in units of the isolated-cell gap $\Delta_{1/2}$ to the lowest excited
doublet, the scale against which the first-order projection is controlled, these
are $g=\beta/\Delta_{1/2}=1.11$, $1.14$, and $1.18$.
They lie outside the range $g\le0.65$ in which finite-chain gaps agree with the
first-order effective chain to within about $5\%$, and at the edge of that
range the continuum upper edge reaches only $55\%$--$59\%$ of the local
quartet gap.
The weak-coupling effective theory therefore predicts a finite energy separation
between the low-energy continuum and the local quartet, not a controlled
hybridization.
For the three local representative values $\alpha=0.14$, 0.23, and 0.33,
raising $\beta$ to 0.5 gives $g=2.24$--$5.32$, well beyond the range where
first-order continuum edges apply.
Figure~\ref{fig:selection} shows the endpoint projection coefficients and the
projector-resolved local weights underlying this comparison of energy scales.

\begin{figure*}[t]
  \includegraphics[width=0.92\textwidth]{fig5_selection_rules.png}
  \caption{Effective-chain projection and local selection rules.
  The geometry and projection coefficients determine the effective exchange.
  Isolated-pentamer levels and momentum-dependent weights identify bright local
  spin-$3/2$ channels and the dark odd-parity doublet at $q=\pi$.}
  \label{fig:selection}
\end{figure*}

\subsection{Collective dynamics}
\label{sec:dynamics}

Figure~\ref{fig:fullq} compares the complete connected momentum response at
$(\alpha,\beta)=(0.5,0.5)$ and $(-0.5,0.5)$ for $L=40$.
Each momentum source is evolved independently using Eq.~\eqref{eq:finite-time-dsf}.
At positive $\alpha$, the low-energy weight is prominent near $q=\pi$ and is
accompanied by higher-energy bands; at negative $\alpha$, low-energy weight is
redistributed among momenta and the higher-energy pattern also changes.
This is a comparison of two finite-chain spectra at fixed resolution.
It does not by itself locate a phase boundary or assign a many-body band to a
particular isolated-cell multiplet.

The relative root-mean-square deviations under $q\to2\pi-q$ are
$2.87\times10^{-6}$ and $4.49\times10^{-6}$ for positive and negative $\alpha$,
respectively, without reflection averaging.
The five-source joint refinement described in Sec.~\ref{sec:dsf} changes the
spectra by relative root-mean-square amounts of $0.79\%$ and $0.45\%$,
respectively, with no shift of the sampled main-peak frequencies.
These checks support the numerical consistency of the displayed spectra at the
tested settings; they do not establish long-time or thermodynamic convergence.

The energy resolution remains $0.09277J_1$.
Independent sector calculations give finite-chain gaps $0.05562J_1$ and
$0.001092J_1$ for the two panels, both smaller than this width.
Consequently, the heatmaps cannot establish that either gap is resolved from
zero frequency, and enhanced weak stripes cannot be identified as additional
branches without further checks of the time window and evolution time.
The isolated-cell parity analysis supplies candidate channels for the
high-energy response; a quantitative correspondence would require
matrix-element-resolved or parameter-continuation evidence in the chain.

\begin{figure*}[t]
  \centering
  \includegraphics[width=0.49\textwidth]{fig6_fullq_positive.pdf}\hfill
  \includegraphics[width=0.49\textwidth]{fig6_fullq_negative.pdf}
  \caption{Complete connected $S_T^{zz}(q,\omega)$ for $L=40$ open chains at
  $J_2/J_1=0.5$ (left) and $-0.5$ (right), both with $J_3/J_1=0.5$.
  The physical-site momentum spans $0\le q\le2\pi$ and energies are in units of
  $J_1$; all 41 momentum sources are computed independently.
  The triangular window at $TJ_1=60$ has width $0.09277J_1$; production settings
  are $\Delta tJ_1=0.05$ and $\chi=192$.
  Both panels use the same blue-to-red color map with normalized color coordinate
  $u(s)=\operatorname{asinh}(s/0.02)/\operatorname{asinh}(6/0.02)$ for
  $s=J_1S_T^{zz}\in[0,6]$.
  This uniform nonlinear display enhances weak weight without changing the data
  or the energy resolution.
  No Gaussian smoothing, interpolation, or reflection averaging is used.}
  \label{fig:fullq}
\end{figure*}

For comparison, Fig.~\ref{fig:cv-tebd} retains an earlier benchmark at
$(\alpha,\beta,L,q)=(0.23,0.5,40,\pi)$ and $\eta=0.1$.
Eleven correction-vector frequencies cover $0\le\omega/J_1\le1.5$ and bracket
the main peak at $\omega=0.05J_1$.
The earlier center-source time-evolution approximation places the peak at the
same frequency; the two estimates differ by at most $6.4\%$ at the shared
frequencies and by $1.0\%$ in their summed sampled weight.
Its real-time protocol uses an 18-site spatial window and differs from the
full-source estimator of Fig.~\ref{fig:fullq}.
The agreement is therefore a restricted benchmark, not a cross-method
validation of the present $\alpha=\pm0.5$ maps or of their high-energy bands.

\begin{figure*}[!tb]
  \includegraphics[width=0.62\textwidth]{fig7_cv_tebd.png}
  \caption{Pointwise correction-vector and time-evolution comparison at
  $\alpha=0.23$, $\beta=0.5$, $L=40$, and $q=\pi$.
  Eleven correction-vector frequencies bracket the principal feature and
  cover $0\le\omega/J_1\le1.5$ at $\eta=0.1$.
  The time-evolution curve uses the earlier center-source, spatial-window
  approximation; it is a separate benchmark from Fig.~\ref{fig:fullq}.}
  \label{fig:cv-tebd}
\end{figure*}

\subsection{Finite-size momentum reconstruction at negative intra-cell exchange}
\label{sec:negative}

The negative-$\alpha$ data show a reproducible change of the dominant
structure-factor momentum.
Deep in the scanned region $\qpeak\simeq0.60\pi$, whereas closer to $\alpha=0$ the
maximum locks to $q=\pi$.

As for positive $\alpha$, the natural first question is what the isolated cell
predicts.
At all 201 sampled points with $-1\le\alpha\le-0.001$, the isolated pentamer has a
spin-$3/2$ ground multiplet rather than the spin-$1/2$ doublet found for
$\alpha>0$.
Projecting the intra-cell couplings onto that quartet and evaluating the
shape factor of Eq.~\eqref{eq:shape} at $q=0.20\pi$, $0.60\pi$, and $\pi$ gives a
positive
effective exchange $\Jeff/\beta=0.108$--$0.111$, with the $q=0.60\pi$ weight in
the range $0.62$--$0.73$, while the $q=\pi$ weight never exceeds
$10^{-2}$ on the same scale.
The projector identities are satisfied to machine precision, and the $0.60\pi$
weight is the largest of the three at every sampled point.
These three physical momenta fold to the same effective-chain antiferromagnetic
momentum $Q=\pi$.
The result therefore selects the $0.60\pi$ replica among these three folded
responses, rather than locating the global maximum of $F(q)$; indeed,
$F(0)=(\sum_a c_a)^2=1$.
The question for the chain is whether its low-energy response retains this
replica weighting and how the dominant momentum moves toward $q=\pi$.

Our operational boundary test seeks an interval in which $\qpeak$ changes and the
finite-size trend of the dominant ratio, $\Delta R=R_{80}-R_{60}$, changes sign
together (Sec.~\ref{sec:diagnostics}).
The chain scan searched for such intervals in two stages.
The first stage evaluated 96 calculations at 48 coupling points; 89
satisfied the convergence criteria of Sec.~\ref{sec:convergence}, while seven
non-boundary cases remained marginal and were excluded.
The converged calculations identified eight candidate intervals across seven
$\beta$ slices.
The second stage then evaluated these intervals with 29 additional $L=60,80$
calculations at 64 sweeps, all of which satisfied the convergence criteria of
Sec.~\ref{sec:convergence}.
Seven of the eight targets retain either a momentum reconstruction or a sign
change in $\Delta R$ larger than twice its uncertainty.
No adjacent interval carries both signatures, only the $\beta=0.05$
momentum interval reaches the target width $0.025$, and the $\beta=0.10$ slice
supplies no supported internal interval at the final size pair.

The data therefore establish a finite-size momentum reconstruction whose two
diagnostics select different intervals. This is compatible with a broad crossover,
a slowly converging incommensurate regime, or residual finite-size effects;
it does not resolve a common thermodynamic boundary.
The converged momentum samples and the two diagnostics are shown in
Fig.~\ref{fig:negative-reconstruction}.

\begin{figure*}[!tb]
  \includegraphics[width=\textwidth]{fig8_negative_reconstruction.png}
  \caption{Negative-$\alpha$ momentum reconstruction from converged calculations.
  Large-size data move from $\qpeak\simeq0.60\pi$ to $\pi$, but the momentum and
  dominant-ratio diagnostics do not select the same adjacent interval.}
  \label{fig:negative-reconstruction}
\end{figure*}

\subsection{Entropy coefficient near \texorpdfstring{$\alpha=-0.10$}{alpha=-0.10}}
\label{sec:saved-state}

Inside the reconstruction, at $(\alpha,\beta)=(-0.10,0.25)$, an earlier $L=80$
analysis gave $\ceff=0.467$--$0.605$, well below the values near unity found at
neighboring couplings.
A genuinely reduced entropy coefficient would point to a distinct regime inside
the crossover, so we tested this point against the initial state, the system
size, and the eigensolver.

Six points were recomputed at $L=80$ with the ground-state wavefunctions saved,
using two $S^z=0$ initial states and 64 sweeps per initial state:
$\alpha=-0.25,-0.175,-0.10$ at $\beta=0.25$ and $\alpha=-0.625,-0.50,-0.425$ at
$\beta=0.75$.
All runs satisfied the convergence criteria of Sec.~\ref{sec:convergence}, with a
largest variance per site of $1.3\times10^{-9}$.
At $\beta=0.25$ the coefficient ranges were $0.467$--$0.605$ at $\alpha=-0.10$,
$0.886$--$0.971$ at $\alpha=-0.175$, and $0.981$--$1.049$ at $\alpha=-0.25$; the
three $\beta=0.75$ points gave ranges within $0.905$--$0.976$.
These ranges measure sensitivity to the fitting window and to the initial state.

A size comparison at $L=40,60,80$ for $\alpha=-0.175,-0.150,-0.125,-0.100$ with
$\beta=0.25$ made the low value look like a trend.
At $\alpha=-0.10$ the coefficients from fits with period-10 bond classes
($b\bmod10$) span
$0.970$--$1.340$, $0.625$--$0.959$, and $0.467$--$0.605$ with increasing length, and
at fixed window $f=0.20$ the mean over the two initial states decreases from
$0.988$ to $0.695$ and $0.533$.
Increasing the Krylov dimension of the eigensolver shows that this decrease is
numerical.
An earlier increase of the Krylov dimension from $k=3$ to $k=8$ revised the
entropy slopes at this point strongly: at $L=80$ the period-10, $f=0.20$
coefficient moves from $0.58$ to $0.93$ for one N\'eel-type initial state and from
$0.48$ to $0.93$ for the other, after which the two initial states agree to
$3\times10^{-8}$.
The further step to $k=12$ changes the observables only by the small amounts
reported in Sec.~\ref{sec:krylov} and leaves $\qpeak=\pi$ unchanged.
The controlled value $\ceff\simeq0.93$ lies within the range $0.886$--$1.049$
found at $\alpha=-0.175$ and $-0.25$, so the suppression of $\ceff$ at
$\alpha=-0.10$ does not survive and no distinct regime is indicated there.
Figure~\ref{fig:negative-size} collects the entropy coefficients, bond-class
variation, connected correlations, and sector gaps behind this assessment.

\begin{figure*}[!tb]
  \includegraphics[width=\textwidth]{fig9_negative_size_evidence.png}
  \caption{Joint finite-size evidence near $\alpha=-0.10$, $\beta=0.25$.
  Entropy coefficients, bond-class variation, connected correlations, and sector
  gaps are shown together to expose size and initialization sensitivity.
  Data use the earlier $k=3$ Krylov setting; the controlled $L=80$, $\alpha=-0.10$ value
  is discussed in the text.}
  \label{fig:negative-size}
\end{figure*}

\section{Discussion}
\label{sec:discussion}

\subsection{Contrast with the tetramer chain}

The main contrast with the tetramer chain follows from unit-cell parity.
Four spin-$1/2$ degrees of freedom can form an integer-spin or singlet cluster
sector, which is why the tetramer problem admits gapped tetramer and Haldane-like
phases separated by a gapless intermediate
region~\cite{li2026tetramer,gong2008tetramer}.
Five spins leave a half-integer degree of freedom unless translation symmetry is
broken or the ground state is degenerate, and the Lieb--Schultz--Mattis-type
constraints~\cite{lieb1961,ogata2019,affleck1987} then exclude a unique
symmetry-preserving gapped ground state under the stated assumptions.
At the tested positive-$\alpha$ points, the surviving doublet converts the
inter-cell exchange into an antiferromagnetic interaction between effective spins.
The finite-size observables accordingly point to a critical effective chain, and
the pentamer has to be classified on its own terms rather than through the
tetramer phase names and their triplon, quarton, or quinton labels.

The spectral comparison sharpens this distinction.
The pentamer's odd-parity excited doublet is dark at $q=\pi$, while higher
even-parity spin-$3/2$ sectors remain bright.
The finite-chain spectra contain low-energy weight and higher-energy structure,
but the local hierarchy alone does not identify the composition of those bands.
The multiplet analysis predicts allowed local channels, whereas a quantitative
assignment in the chain requires additional matrix-element or continuation
calculations.
A qualitatively similar separation between momentum-selective local weight and a
collective continuum has been discussed for other cluster and chain
magnets~\cite{keselman2020,gao2024srv3o9,bruognolo2016}, in line with the
interpretation given here.

The odd-cell constraint excludes a unique symmetric gapped ground state while
leaving open gapped symmetry-broken phases, degenerate ground states, and gapless
phases of other origin.
Among these, the finite-size evidence at the tested points selects the $c=1$
effective chain.
String-order and entanglement-spectrum
diagnostics~\cite{guwen2009,pollmann2010,pollmann2012,han2017bcfg} would be the
natural tools for probing hidden order in the same framework.

\subsection{Controlled interpretation near \texorpdfstring{$\alpha=0$}{alpha=0}}

Both sides of the coupling plane meet at $\alpha=0$, where the end bonds of each
cell vanish.
Sites 2--4 then form a trimer with a spin-$1/2$ ground doublet, while site 5 of
one cell and site 1 of the next are joined only by $J_3$ into a singlet dimer.
For $|\alpha|\ll\min(1,\beta)$ the weak end bonds let neighboring trimer doublets
interact through a virtual triplet excitation of the dimer, which costs an energy
of order $\beta$.
Once the end spins are bound into these dimers, the trimer doublets are the
low-energy degrees of freedom for either sign of $\alpha$, even though the isolated
cell, which lacks the inter-cell bond, has a quartet ground state for $\alpha<0$.
Projecting onto the trimer doublets and eliminating the dimer triplet gives
\begin{equation}
\Jeff=\frac{2\alpha^2}{9\beta}>0 .
\label{eq:near-zero}
\end{equation}
Exact diagonalization of an eight-spin trimer--dimer--trimer cluster at 47 points
supports this
local $\alpha^2/\beta$ scaling in its asymptotic domain, consistent with the
effective-spin logic established for trimer
chains~\cite{cheng2022trimer,cheng2024trimerfield}.
The finite-$\alpha$ controls at $(-0.10,0.25)$ and $(-0.175,0.25)$ deviate from the
leading term by $-36.1\%$ and $-53.9\%$, respectively.
These are signed deviations of the local approximation itself, defined as
$\epsilon=(\Delta_8/\alpha^2-2/(9\beta))/(2/(9\beta))$, where $\Delta_8$ is the
singlet--triplet gap of the eight-spin cluster, which for two coupled effective
spins equals $\Jeff$; the deviations grow as
$\alpha$ leaves the asymptotic domain, as expected for a leading-order result.
The decomposition thus supplies the mechanism near $\alpha=0$: it explains why the
effective exchange is antiferromagnetic and why it scales as $\alpha^2/\beta$.
It also explains why the near-zero strip is the most delicate region numerically:
the many-body energy scale vanishes quadratically as $\alpha\to0$, so finite-size
and solver effects compete with a very small exchange there.

\subsection{Meaning and boundary of the negative-\texorpdfstring{$\alpha$}{alpha} result}
\label{sec:negative-boundary}

The negative-$\alpha$ scan finds a change of the dominant finite-size momentum
across the sampled coupling plane at the retained solver settings.
In the thermodynamic limit such a reconstruction could sharpen into a Lifshitz
transition, spread into a broad incommensurate regime, or remain a crossover
driven by the competing-momentum mechanism known from frustrated and dimerized
chains~\cite{hida1992,chitra1995,hakobyan2001,white2008s1}.
The mismatch between the two diagnostics is compatible with a crossover but does
not distinguish these possibilities in the thermodynamic limit.
The local quartet analysis supplies a candidate origin for the deep-region value:
among the three inequivalent momenta folded from $Q=\pi$, the isolated
spin-$3/2$ cell has its largest shape-factor weight at $0.60\pi$ throughout the
sampled negative region, matching the measured deep-region $\qpeak$.
This is a comparison of folded replicas, not a claim that the isolated shape
factor has its global maximum there.
Whether the extended chain develops a singularity at that replica is a separate
many-body question, and the near-zero strip is especially sensitive to
initialization and fitting window.
The correspondence motivates a test: in the regime where the quartet projection
is controlled, larger-chain spectra should approach the relative replica weights
of Eq.~\eqref{eq:folded-response}.

The refinement calculations pass the stated acceptance gates, and the
$L=80$, $(\alpha,\beta)=(-0.10,0.25)$ endpoint is stable under
$k=8\rightarrow12$.
That local control does not exclude solver bias throughout the negative-region
scan; in particular, the $L=40$ and 60 members of the entropy size sequence still
use $k=3$.
Momentum reconstruction and finite-size flow are distinct observables, and at the
available sizes and coupling resolution they select different intervals.
The present evidence supports a finite-size momentum reconstruction.
A thermodynamic phase assignment requires a consistently controlled size sequence
and an appropriate bulk diagnostic; coincident signatures in the present
operational test would not alone establish such an assignment.

\subsection{Outlook}
\label{sec:remaining-tests}

The negative-$\alpha$ result sets up a sharp test of the phase question.
An observable with a controlled large-size singularity or crossing that tracks the
momentum reconstruction would test a thermodynamic boundary.
Figure~\ref{fig:fullq} now provides a negative-$\alpha$ dynamical example at
$(-0.5,0.5)$; a coupling scan and longer-time, larger-size spectra are needed to
connect that example to the static reconstruction intervals.
The five-momentum joint refinement should also be extended across the full grid,
with separate time-step and bond-dimension checks where quantitative line shapes
are required.

Two further extensions would connect the present results to the wider literature
on one-dimensional spin dynamics.
First, the momentum-resolved line shapes could be compared with the quasiparticle
interaction signatures and spinon hydrodynamics that have been analysed in
related chains~\cite{keselman2020,balents2020,wang2022hydro}, and with the
dynamical signatures used to distinguish symmetry-broken from liquid
phases~\cite{drescher2023,li2023dqcpsteer}.
Second, the entanglement analysis could be extended beyond the single-scalar
$c_{\mathrm{eff}}$ reported here to the full entanglement spectrum and its
dependence on microscopic anisotropy~\cite{kargarian2008}, and the momentum-space
continuum could be analysed with the same tools that have been applied to
honeycomb and square-lattice
antiferromagnets~\cite{gu2022honeycomb,wu2022j1j3}.

\subsection{Relation to experiment}

The local selection rule has a direct experimental consequence.
For an isolated reflection-symmetric cell, parity removes the lowest excited
doublet from the $q=\pi$ single-spin channel.
In a sufficiently weakly coupled pentamer material, this motivates a suppressed
lowest-doublet signal and stronger allowed spin-$3/2$ channels in inelastic
neutron scattering or single-spin $L$-edge RIXS.
Coupling-induced mixing and the relevant scattering operator must be checked
before extending the local selection rule quantitatively to a material.
Such a comparison follows the spectral hierarchy used to
identify composite excitations in trimer and tetramer
compounds~\cite{cheng2022trimer,li2025rixstrimer,li2026tetramer,prabhakar2025trimer}.
The prediction concerns the single-spin response, which is the channel measured
by neutrons and $L$-edge RIXS.
Because the Hamiltonian is SU(2) symmetric and the finite-chain ground state at
the calculated points is a spin singlet, $S^{zz}(q,\omega)$ fixes the unpolarized
neutron response up to a factor of three.
Two-particle channels probed by $K$-edge RIXS, and their relation to the
multi-spinon continuum observed in one-dimensional
cuprates~\cite{schlappa2018,ishii2025kedge,kumar2020,forte2011}, are a natural
extension of the same framework.

\section{Conclusion}
\label{sec:conclusion}

We have traced the behavior of a spin-$1/2$ Heisenberg pentamer chain from
isolated local multiplets to collective response.
At the tested positive couplings, the pentamer retains an even-parity spin-$1/2$
doublet, positive inter-cell exchange generates an antiferromagnetic effective
chain, and the leading projection error is quadratic in the inter-cell coupling.
Entanglement, finite-size gaps, correlations, and static structure factors
identify a $c=1$ effective chain at those points.
Reflection parity makes the lowest excited doublet dark at $q=\pi$ and shifts most
local inelastic weight into spin-$3/2$ sectors.
Complete connected real-time spectra at $(\alpha,\beta)=(\pm0.5,0.5)$ show
different distributions of low- and high-energy weight in open $L=40$ chains.
Their common finite-time width is $0.09277J_1$, and a joint refinement supports
stability at five tested momenta.
The earlier correction-vector comparison is restricted to a different parameter
point and real-time estimator; the collective bands have not yet been assigned
quantitatively to individual local multiplets.

On the negative-$\alpha$ side, the isolated cell switches to a spin-$3/2$ quartet
whose shape factor favors the $0.60\pi$ replica among folded $Q=\pi$ momenta
at every sampled point, and the
accepted calculations reveal a reconstruction from $\qpeak\simeq0.60\pi$ to
$q=\pi$.
All 29 final refinement calculations converge, and a
terminal $k=8\rightarrow12$ control shows that the most sensitive $L=80$ endpoint
is stable with respect to a larger Krylov subspace.
This endpoint control does not establish solver stability of the entire size
sequence or scan.
None of the eight candidate intervals satisfies the momentum and dominant-ratio
criteria simultaneously, so the reconstruction remains a finite-size result
compatible with a crossover rather than an established thermodynamic boundary.

For an odd-unit-cell cluster chain, the low-energy physics is controlled not by an
inherited phase label but by which local multiplet survives, how it couples
between cells, and which of its excitations the symmetry of the probe allows one
to see.


\begin{acknowledgments}
[Add funding, computing resources, institutional support, and individual
acknowledgments.]
\end{acknowledgments}

\section*{Data and code availability}
The Hamiltonian construction, the DMRG and spectral drivers, the figure-generation
scripts, and the numerical data underlying
Figs.~\ref{fig:multiplets}--\ref{fig:negative-size} are available at
[repository URL or archival DOI, release identifier].

% 压缩参考文献条目间距。
\setlength{\bibsep}{2.0pt plus 0.3ex}
\bibliography{references}

\end{document}
